\documentclass[12pt,a4paper,draft]{report}
\usepackage[left=1cm, right=1cm, top=2cm, bottom=2cm]{geometry}
\usepackage{lscape} % Paket für Querformat
\usepackage{booktabs} % Für bessere Tabellenlinien
\usepackage[ngerman]{babel}
\usepackage[T1]{fontenc}
\usepackage{colortbl} % Für Hintergrundfarben in Tabellen
\usepackage{array} %für bessere Tabellenformatierung
\usepackage{xcolor}
\usepackage{graphicx}
\usepackage{tabularx}
\usepackage{babel}
\usepackage{chemformula} % Hauptpaket für chemische Formeln
\title{Tabellen für Zwischenpräsi}
\author{Linus Benetz}

\begin{document}
	\section{Charecteristics M. thau und M. bovis}
\begin{table}[h]
	\centering
	\begin{tabularx}{\linewidth}{l c >{\centering\arraybackslash}X}
		\toprule
		\rowcolor[gray]{0.9} & \textit{Methanobrevibacter thaueri} CW\textsuperscript{T}& \textit{Methanobrevibacter boviskoreani} AbM4\\
		\midrule
		form & coccobacillus & rods \\
		
		G + C content (mol\%) &38&29\\
		
		T-Optimum & 37 °C & 38 °C \\
		
		pH-Optimum &7 & 6.5-7\\
		nutritional requirements & acetat and/or yeastextract/trypticase & yeastextraxt\\
		
		growth on \ch{EtOH} & $-$& $+$\\
		
		growth on formate & $-$& $+$\\
		
		antibiotic sensitiv & ND & chloramphenicol, puromycin *\\
		
		antibiotic tolerance & ND & bacitracin, vancomycin *\\
		\bottomrule
		*in \textit{M. boviskoreani JH1}\textsuperscript{T}&&
		
		
		
		
		
		
	\end{tabularx}
\end{table}


\section{AB with organism}

\begin{table}[h]
	\begin{tabularx}{\linewidth}{l c *{2}{>{\centering\arraybackslash}X}}
		\toprule
		\rowcolor[gray]{0.9} antibiotic & MIC [µg/ml] & organism  & paper\\
		\midrule
		\rowcolor{blue!5} Puromycin & 2 & \textit{Methanococcus voltae} & Possot et al 1988  \\
		
		\rowcolor{blue!5} & 5 & \textit{Methanosarcina mazei} & Mondorf et al 2012 \\ 
		
		Acetylpuromycin & 100 & \textit{Methanococcus voltae} &  Possot et al 1988 \\ 
		
		\rowcolor{blue!5}Neomycin & 20 & \textit{Methanosarcina mazei} & Mondorf et al 2012\\
		
		\rowcolor{blue!5}	& 125&\textit{Methanobrevibacter arboriphilus}& Pecher \& Böck 1988 \\ 
		
		\rowcolor{blue!5}	& 100 &\textit{Methanothermobacter thermautotrophicus}& Fink et al. 2023 \\ 
		
		Clarithromycin & 0.25 - 16& \textit{Mycobacterium abscessus}&  Liao et al 2024  \\
		
		\rowcolor{blue!5} Chloramphenicol & 25 & \textit{Methanobrevibacter smithii}& Dridi et al. 2011 \\
		
		\rowcolor{blue!5} &4& \textit{Methanosphaera stadtmanae}& Dridi et al. 2011\\ 
		
		\rowcolor{blue!5} Thiamphenicol  & ND & & \\ 
		
		Bacitracin &4&\textit{Methanobrevibacter. smithii}& Dridi et al. 2011 \\
		\bottomrule
	\end{tabularx}
\end{table}	
	\section{AB with resistance gen}
	
	\begin{table}[h]
		\begin{tabularx}{\linewidth}{l c *{2}{>{\centering\arraybackslash}X}}
			\toprule
			\rowcolor[gray]{0.9} antibiotic & MIC [µg/ml]* & effects  & resistance gene\\
			\midrule
			\rowcolor{blue!5} Puromycin & 2;5 & AA-tRNA-analoga -> attacks protein-biosynthesis &  \textit{pac} \\ 
			
			Acetylpuromycin & 100 &  & product of \textit{pac}  \\ 
			
			\rowcolor{blue!5}Neomycin & 20;100;125 & binds to 16s-Subunit -> misreading of mRNA & \textit{neoR}\\
			
			Clarithromycin & 0.25;16& binds to 50s Subunit -> prevents transfer of amino acid chain &  \textit{ermB}  \\
			
			\rowcolor{blue!5} Chloramphenicol & 4;25 & binds to peptidyltransferase of 50s-Subunit -> prevents protein biosynthesis& \textit{catP} \\
			
			\rowcolor{blue!5} Thiamphenicol  & ND & analoga of Chloramphenicol & \textit{catP}\\ 
			
			Bacitracin &4&inhibits transport of cellwall building-blocks& \textit{BcrABC}; \textit{BacRS}; \textit{BacA}  \\
			\bottomrule
			*tested in close relatives&&&
		\end{tabularx}
	\end{table}
% Mondorf et al 2012 and Possot et al 1988 and Pecher \& Böck 1988 Liao et al 2024 
% & Dridi et al. 2011&

\end{document}